\section{总结与延伸}

\subsection{关键要点}
\begin{itemize}
    \item ``Dive into Kimi K2.5'' 把整份讲稿框在 data、algorithm、infra 三条主线里，方便我们按逻辑拆解技术报告。
    \item 早期就混入视觉数据、用 benchmark 驱动的合成 pipeline 是 K2.5 能在图文、代码、GUI 上打通能力的基础。
    \item GRM 成为统一 reward，Pretrain/SFT/RL 三段训练叠加出稳态能力，deg-and-recover 是正常的收敛轨迹。
    \item Agent swarm 的魅力在于可学的 Allcastrater、并行 subagent 与 critical steps 资源约束，使得多 Agent 比单 Agent 更高效。
    \item Infrastructure 通过 Critical Steps + Kimi code bench 进行评估，即便讲师少解释，训练曲线和 bandmark 已经证明了该体系可行。
    \item 实战部分展示了如何把字幕、关键帧、slides 结合成具备 summary table 与 box 层次的文档，确保 audit 通过。
    \item 最终，原生多模态 + Agent Swarm 让 K2.5 真正具备了 ``Visual Agency'' 的执行力，而不仅是视觉理解。
\end{itemize}

\subsection{总结表}
\begin{table}[H]
\footnotesize
\centering
\begin{tabular}{p{3cm}p{7cm}p{4cm}}
\toprule
维度 & 洞见 & 证据 \\
\midrule
Data & 早期混入视觉、定义 benchmark，再用合成 pipeline 逐阶段扩展能力。 & 附录 Figure 9 中 1:9~1:1 的比例实验；讲师列出的三步数据构建流程。 \\
Algorithm & GRM + PARL RL 统一 reward，多阶段训练让 open-ended 任务有稳定度；Deg \& Recover 是自然波动。 & GRM 知识盒与 warningbox 里的解释，PARL agent loop、critical steps 训练曲线。 \\
Agent \& Infra & Allcastrater 确保 subagent 调度、Critical Steps 限定资源，Evaluation 使用 Kimi code bench 与 bandmark。 & Agent section 中 createSubagents/assignTask 讲解 + infra section 提到评估 bandmark。 \\
Practice & 实战章节规范化素材处理、关键帧与 slides 的插图、summary table，确保文档满足 audit 要求。 & 实战章节中的 knowledgebox、importantbox 以及 figure timecode 示例。 \\
\bottomrule
\end{tabular}
\caption{从数据、算法到 agent/infra 的层层产业化，体现 K2.5 原生多模态 + Agent Swarm 架构的逻辑闭环。}
\end{table}

\subsection{拓展阅读}
\begin{itemize}
    \item \href{https://www.kimi.com/ai-models/kimi-k2-5/}{Kimi K2.5 | Open Visual Agentic Model for Real Work} — 官方模型概览与接口说明。
    \item \href{https://www.kimi.com/blog/kimi-k2-5}{Kimi K2.5 Tech Blog: Visual Agentic Intelligence} — 详细介绍 Visual Agency Intelligence 与数据策略。
    \item \href{https://artificialanalysis.ai/articles/kimi-k2-5-everything-you-need-to-know}{Kimi K2.5 — Everything you need to know (Artificial Analysis)} — 外部视角的 benchmark 与 agent swarm 小结。
    \item \href{https://futuretechdiaries.com/article.php?slug=kimi-k2-5-a-deep-dive-into-chinas-long-context-ai-model}{Kimi K2.5: A Deep Dive into China’s Long-Context AI Model} — 比较 + 适用场景分析。
    \item \href{https://platform.kimi.ai/docs/guide}{Kimi API Platform Guide} — 说明视觉 agent 调用和工具链的实践方法。
\end{itemize}

\subsection{本章小结}
这场串讲最终告诉我们：K2.5 构建了一个原生多模态的 agentic stack，数据、GRM 驱动的算法与 infra/agent 协调形成闭环，赋予模型真正的 Visual Agency Intelligence。
%% --- Fin del contenido --- %%

\end{document}
